
\subsubsection{Pre-Training Data Curation}

Large-scale pre-training datasets apply deduplication, perplexity filtering, and language detection to web-scraped corpora. C4 \citep{Raffel2020Exploring} documents filtering with deduplication threshold ~0.8 and perplexity cutoff ~1000. The Pile \citep{Gao2020ThePile} and RedPajama curate diverse text sources. DataComp \citep{Gadre2023DataComp} systematically explores vision-language pre-training filters, demonstrating that curation quality significantly impacts downstream performance. These works establish stage-specific best practices but do not test whether discovered thresholds transfer to fine-tuning or RLHF stages. Our work extends this line by characterizing which pre-training curation decisions can be reused downstream.

\subsubsection{Fine-Tuning Data Selection}

Instruction tuning research emphasizes quality over quantity. LIMA \citep{Zhou2023LIMA} achieves strong performance with 1,000 high-quality examples. Alpagasus \citep{Chen2023Alpagasus} filters Alpaca-52k using ChatGPT-based quality scoring. These methods optimize for instruction-following tasks without leveraging curation knowledge from pre-training. Our taxonomy shows that task-specific quality filters require stage-specific tuning (5.04\% transfer degradation), while low-level hygiene operations transfer robustly from pre-training (0.38\% delta).

\subsubsection{RLHF and Preference Data Curation}

RLHF pipelines curate preference data to align models with human values (Ouyang et al., 2022; Bai et al., 2022). Constitutional AI \citep{Bai2022Constitutional} generates synthetic preference pairs, while Anthropic HH-RLHF manually curates helpfulness and harmlessness examples. This work treats RLHF curation independently from earlier stages. Our framework predicts that safety filters (universal hygiene) should transfer from pre-training, while preference alignment criteria (objective-dependent) require RLHF-specific tuning. This remains an untested hypothesis for future work.

\subsubsection{Transfer Learning in Foundation Models}

Transfer learning literature studies how model parameters transfer across tasks (Howard and Ruder, 2018; Devlin et al., 2019; Brown et al., 2020). Pre-trained representations generalize to downstream tasks with minimal fine-tuning. Our work applies transfer analysis to data curation rather than model parameters, examining which curation decisions exhibit cross-task robustness. We find that data hygiene operations (deduplication, outlier removal) transfer similarly to how low-level features transfer in model representations: both address universal properties independent of specific objectives.

\subsubsection{Subset Selection and Coreset Methods}

Coreset construction \citep{Feldman2011Scalable} and k-center greedy subset selection \citep{Sener2018Active} reduce dataset size while preserving representativeness. DataComp applies embedding-based diversity maximization for pre-training. Our h-m3 experiment extends this by quantifying the quality-speed trade-off when embedding stage mismatches curation stage: early-stage embeddings (MiniLM) are 6.$9\times$ faster than late-stage (Instructor) but incur 2.4\% quality penalty.

\subsubsection{Positioning}

Existing work optimizes curation per-stage without testing cross-stage transfer. We provide the first systematic characterization of which curation techniques transfer robustly (low-level hygiene) versus which require stage-specific tuning (high-level strategy). This taxonomy enables practitioners to reuse universal operations while focusing optimization on stage-dependent components.
